% !TEX root = chapter-standalone.tex

\section[Dynamical systems]{Dynamical systems}
\label{sec:cat-dyn}


\

\begin{ctdefinition}[Endofunctions]
    \label{def:cat-endofunctions}
    The category $\EndSet$ of endofunctions is:
    \begin{enumerate}
        \item \emph{Objects}: pairs $\tup{\setA, \stylemorph{\varphi}}$ where $\setA$ is a set and $\stylemorph{\varphi} \colon \setA \to \setA$ is a function.
        \item \emph{Morphisms}: a morphism $\mora \colon \tup{\setA, \stylemorph{\varphi}} \mto_\EndSet \tup{\setB, \stylemorph{\psi}}$ is a function $\mora \colon \setA \to_\Set \setB$ with the property that $\mora \mthen \stylemorph{\psi}= \stylemorph{\varphi} \mthen \mora$.
        \item \emph{Composition}: the usual composition of functions.
        \item \emph{Identity morphisms}: identity functions.
    \end{enumerate}
\end{ctdefinition}

\begin{exercise}
    Prove that \cref{def:cat-endofunctions} is indeed a category.
\end{exercise}

\begin{solution}
    \todotext{Write this solution}
\end{solution}

\begin{gradedexercise}[\exname{Endofunctions}]
    \label{ex:endofunctions}

    In this exercise we will define a category $\EndSet$ of ``endofunctions'' and it is your task to check that it really is a category.

    The objects of $\EndSet$ are pairs $\tup{\setA, \stylemorph{\varphi}}$ where $\setA$ is a set and $\stylemorph{\varphi} \colon \setA \to \setA$ is a function.

    A morphism $\mora \colon \tup{\setA, \stylemorph{\varphi}} \mto_\EndSet \tup{\setB, \stylemorph{\psi}}$ is a function $\mora \colon \setA \to_\Set \setB$ with the property that $\mora \mthen \stylemorph{\psi}= \stylemorph{\varphi} \mthen \mora$.

    Composition is given by composition of functions, and identity morphisms are identity functions.

\end{gradedexercise}

\solutionof{Endofunctions}
